% ANCHOR: after_label def:c10-letters
\begin{remark}[Physics notation: letters and composite operators]
A letter is one constituent field of a composite operator with its derivatives, in the components of the physics notation key after Section 1.2: $\partial^\alpha c$, $\partial^\alpha A_m$, $\partial^\alpha A_I$ (equivalently $\partial^\alpha A_{ij}$ in the antisymmetric $\mathbf 6$ of $SU(4)$), $\partial^\alpha\lambda_\alpha{}^i$ (undotted, $S_+$, in the $\mathbf 4$), $\partial^\alpha\bar\lambda_{\dot\alpha i}$ (dotted, $S_-$, in the $\bar{\mathbf 4}$), and the antifields $\partial^\alpha A^*_m$, $\partial^\alpha A^*_{ij}$, $\partial^\alpha\lambda^{*\alpha}{}_i$, $\partial^\alpha\bar\lambda^{*\dot\alpha i}$, $\partial^\alpha c^*$. Its Weyl weight is its canonical mass dimension (block after Definition \ref{def:c6-sigma}) and its degree is its ghost number (key). A gauge-invariant operator is a polynomial in the covariant letters $D^nF_{mn}$, $D^nA_I$, $D^n\lambda$, $D^n\bar\lambda$, the letters of [B03c, Table 1]; written in the letters of this chapter, each covariant derivative $D_m=\partial_m+[A_m,-]$ contributes letters $A_m$ with fewer derivatives, so a covariant letter is a sum of monomials of ordinary letters. The pair weight of two letters is the sum of their canonical dimensions.
\end{remark}

% ANCHOR: after_label thm:c10-decomp
\begin{remark}[Physics notation: one-loop renormalisation acts on pairs of fields]
Theorem \ref{thm:c10-decomp} is the statement that the one-loop divergence of a composite operator is the sum over the pairs of its constituent fields of the Wick contraction of the pair with the propagator in the background of all fields, the propagator dressed by any number of insertions of the interaction vertices; the remaining fields are spectators. This is the two-body form of the one-loop dilatation operator, which acts on two letters at a time [B03c, (2.12), (2.16)]. The term $k=0$ is the self-contraction of the two fields, a pure power divergence (the normal-ordering term), and the graphs in which the operator sits on a tree are the field-strength and coupling renormalisation of the action (Chapter 4), which carry no counterterm of the operator. The graphs of the scalar bilinears are those of the block after Proposition \ref{prop:c5-census}.
\end{remark}

% ANCHOR: after_label lem:c10-colour
\begin{remark}[Physics notation: colour factors and the planar limit]
For $\mathfrak{su}(N)$ with $(X,Y)=-\tfrac{4\hbar}{g^2}\operatorname{tr}(XY)$, $\kappa=-\tfrac{4\hbar}{g^2}$, and with the fields $X=-igX^aT^a$ of the key, $(-,-)^{-1}=2\kappa^{-1}\sum_aT^a\otimes T^a$ for $\operatorname{tr}(T^aT^b)=\tfrac12\delta^{ab}$, and the rules of Lemma \ref{lem:c10-colour} (ii) are the familiar contraction identity $\sum_a(T^a)_{ij}(T^a)_{kl}=\tfrac12\big(\delta_{il}\delta_{jk}-\tfrac1N\delta_{ij}\delta_{kl}\big)$ applied to the propagator of the pair. The term with $N\operatorname{tr}(M_1Y_1\cdots Y_r)$ is the planar contraction of two neighbouring letters of a single trace, of order $g^2N$ in the unit $-\frac{g^2N}{32\pi^2}$; the other terms split the trace or are suppressed by $\tfrac1N$, the non-planar part of the one-loop dilatation operator. In terms of the 't Hooft coupling $g^2N$ the unit is $-\frac{g^2N}{32\pi^2}$.
\end{remark}

% ANCHOR: after_label prop:c10-bound
\begin{remark}[Physics notation: power divergences and operator mixing]
The proper-time cutoff $\varepsilon$ has dimension length$^2$, and $\varepsilon^{-1/2}$ is the momentum cutoff (block after Section 6.1). The coefficient of $\varepsilon^{-m}$ is a power divergence $\propto\varepsilon^{-m}$, a mixing of the pair into composite operators of dimension lower by $2m$, and the coefficient of $\log\varepsilon$ a logarithmic divergence, a mixing into operators of the same dimension, the source of the anomalous dimensions. The bound on the number of output fields is dimensional analysis together with ghost-number conservation: a pair of total dimension $w$ and ghost number $d$ mixes into at most $w+d$ fields; for example $(\partial_m\lambda_\alpha{}^i,\bar\lambda_{\dot\alpha j})$, of dimension $4$ and ghost number $0$, mixes into at most four fields, with power divergences up to $\varepsilon^{-2}$. In the canonical normalisation of the key ($A_m=-igA^a_mT^a$, $A_I=-igA^a_IT^a$), a term of a pair kernel with $r$ output fields multiplies the operator by $-\frac{g^2N}{32\pi^2}(-ig)^{r-2}$ times a number: order $g^2$ for two outputs and $g^3$ for three, the length-changing terms of the spin-chain picture. At two loops at most three fields are contracted (Remark \ref{rem:c10-boundread}).
\end{remark}

% ANCHOR: after_label thm:c10-brst
\begin{remark}[Physics notation: the Slavnov--Taylor identity of operator renormalisation]
$Q+\{I,-\}$ is the BRST variation $s$ on functionals of the fields and the BV differential with the BRST sources $A^*_m$, $A^*_{ij}$, $\lambda^{*\alpha}{}_i$, $\bar\lambda^{*\dot\alpha i}$, $c^*$ (block after Section 6.3). Theorem \ref{thm:c10-brst} states that the logarithmic counterterm of the BRST variation of an operator is the BRST variation of its logarithmic counterterm, up to the field-strength renormalisation of the constituent fields (the block after Lemma \ref{lem:c6-count}): the one-loop Slavnov--Taylor identity of composite-operator renormalisation. Twice the map of the theorem is the one-loop anomalous-dimension operator in the Callan--Symanzik sense of Theorem \ref{thm:c6-dil} (block after that theorem): twice the logarithmic operator counterterm plus the wave-function renormalisation $\gamma_{A_I}=\gamma_{A_m}=\frac{g^2N}{8\pi^2}$, $\gamma_\lambda=\frac{g^2N}{4\pi^2}$, $\gamma_c=\frac{g^2N}{16\pi^2}$ of the letters (block after Section 6.4). (W1) is the statement for one field, (W2) for two, and Proposition \ref{prop:c10-reduction} says that the BRST identity completes the outputs into covariant combinations, $\partial_m\to D_m=\partial_m+[A_m,-]$ and $\partial_mA_n-\partial_nA_m\to F_{mn}$, once the part with fewer fields is known.
\end{remark}

% ANCHOR: after_label thm:c10-simplex
\begin{remark}[Physics notation: Feynman parameters and the regularisation scheme]
$\alpha_j=t_j/\sum_it_i$ are the Feynman parameters of the one-loop graph, the centre $-\sum_j\alpha_jq_j$ is the shift of the loop momentum that completes the square, and the momentum form $\sum_{i<j}\alpha_i\alpha_j|q_i-q_j|^2$ is the second Symanzik polynomial of the one-loop graph, the first being $\sum_j\alpha_j=1$. The regulator cuts every Schwinger proper time at $t_e\ge\varepsilon$, which in Feynman parameters is the cutoff $\tau\ge\varepsilon/\min_e\alpha_e$ on the total proper time; Proposition \ref{prop:c10-shape} states that the logarithmic divergence, and with it every anomalous dimension, is the same for any cutoff shape, while the power divergences and the finite parts are scheme dependent. All coefficients are rational: Beta-type integrals of polynomials in the Feynman parameters.
\end{remark}

% ANCHOR: after_label prop:c10-bubble
\begin{remark}[Physics notation: point splitting and the distribution of derivatives]
The generating function places the two fields at $x+y_a$ and $x+y_b$; the coefficient of $y_a^\alpha y_b^\beta/\alpha!\beta!$ is the kernel of the pair $(\partial^\alpha\varphi_a,\partial^\beta\varphi_b)$, so the derivatives of the two letters are redistributed over the outputs by the Feynman-parameter integrals $\int_0^1\alpha^{n_1+j}(1-\alpha)^{n_2+j}d\alpha=\frac{(n_1+j)!\,(n_2+j)!}{(n_1+n_2+2j+1)!}$. In Beisert's harmonic action [B03c, (4.27)--(4.29)] the two-site Hamiltonian redistributes the oscillators of two letters with the coefficients $(-1)^{1+n_{12}n_{21}}\Gamma(\tfrac12(n_{12}+n_{21}))\Gamma(1+\tfrac12(n-n_{12}-n_{21}))/\Gamma(1+\tfrac12n)$, also Beta functions, with the diagonal value $h(\tfrac n2)$; on the $\mathfrak{sl}(2)$ sector it reduces to $H'_{12}$ of [B03c, (3.14)], which Proposition \ref{prop:c10-twist} reproduces on two sites from these integrals summed over the chains. For the quartic bubble, a pair $(\partial_\mu\phi,\partial_\nu\phi)$ of scalars is replaced by $\frac{\hbar}{16\pi^2}\log\varepsilon\big(-\tfrac16\partial_\mu\partial_\nu-\tfrac1{12}\delta_{\mu\nu}\partial^2\big)$ of the product of the two outputs (Example \ref{ex:c10-closed}).
\end{remark}

% ANCHOR: after_label prop:c10-W0
\begin{remark}[Physics notation: the bosonic pair kernels in components]
In the tables, $a_1$, $a_2$ are the polarisations of the two letters and $b_j$ those of the outputs: $(a_1,b_1)=a_{1m}b_{1m}$ for gluons and $\tfrac14a_1^{ij}b_{1ij}=a_{1I}b_{1I}$ for scalars, and $ip_j=\partial$ on the $j$-th output field. For example the rows of $(A_I,A_J)$ with the outputs $(A_K,A_L)$ are the gluon-exchange triangle (two cubic vertices joined by a gluon line, coefficient $2\delta_{IK}\delta_{JL}$) and the quartic scalar vertex ($-2\delta_{IJ}\delta_{KL}-2\delta_{IK}\delta_{JL}+4\delta_{IL}\delta_{JK}$), each times the colour factor of the table; their sum gives the $SO(6)$ density of Proposition \ref{prop:c10-so6}. Remark \ref{rem:c10-twocomp} gives the remaining two-component forms.
\end{remark}

% ANCHOR: after_label prop:c10-so6
\begin{remark}[Physics notation: the $SO(6)$ spin chain]
With $\gamma=2\times$(the map) the anomalous dimension (Theorem \ref{thm:c10-brst}) and the unit $-\frac{g^2N}{32\pi^2}$, Proposition \ref{prop:c10-so6} is the one-loop dilatation operator of the scalar sector,
\[
\gamma=\frac{g^2N}{8\pi^2}\sum_{l}\big(1-P_{l,l+1}+\tfrac12K_{l,l+1}\big),
\]
the $SO(6)$ integrable spin chain of the scalar sector, with $P$ the permutation and $K$ the trace of the two neighbouring $SO(6)$ indices; in the $SU(4)$ form, $(A_{ij},A_{kl})\mapsto2A_{kl}A_{ij}-\tfrac12\varepsilon_{ijkl}A^{mn}A_{mn}$ in units. It gives $0$ on the $\mathbf{20'}$ and $\frac{3g^2N}{4\pi^2}$ on the Konishi operator (block after Remark \ref{rem:c6-konishi}); on the $\mathfrak{su}(2)$ sector of Proposition \ref{prop:c10-su2}, with $Z=A_1+iA_2$, $X=A_3+iA_4$, it is $\frac{g^2N}{8\pi^2}$ times the Heisenberg Hamiltonian $\sum_l(1-P_{l,l+1})$, whose cyclic spectra are those of the table.
\end{remark}

% ANCHOR: after_label prop:c10-twist
\begin{remark}[Physics notation: twist two and harmonic sums]
$n=(1,i,0,0)$ is a null vector of the complexified Euclidean space, the Wick image of a light-cone direction, and $(\iota_nd_{A|_V})^kZ=(n^mD_m)^kZ$ are the light-cone covariant derivatives of the complex scalar $Z=A_1+iA_2$; the operators $\operatorname{tr}\big((n\cdot D)^kZ\,(n\cdot D)^{S-k}Z\big)$ are the twist-two operators of the $\mathfrak{sl}(2)$ sector. The eigenvalues $-8h(S)$ of the map give the one-loop anomalous dimensions
\[
\gamma(S)=\frac{g^2N}{2\pi^2}\sum_{j=1}^S\frac1j\qquad(S\ \text{even}),
\]
the harmonic sums of [B03c, (2.31)]: $0$ for the $\mathbf{20'}$ descendants at $S=0$ and $\frac{3g^2N}{4\pi^2}$ at $S=2$, the Konishi multiplet. Twist two needs the gluons inside the covariant derivatives: without the pairs $(A_m,Z)$ the value at $S=2$ would be $\frac{g^2N}{6\pi^2}$ instead of $\frac{3g^2N}{4\pi^2}$.
\end{remark}

% ANCHOR: after_label cor:c10-beisert
\begin{remark}[Physics notation: comparison with the one-loop dilatation operator]
[B03c, (1.2), (1.3)] writes the dilatation generator as $D(g)=D_0+\frac{g^2N}{8\pi^2}H+O(g^3)$, $H=\sum_lH_{l,l+1}$ with $H_{12}=2h(J_{12})$ the harmonic action on two neighbouring letters. With the physics dilatation $D=-E=-\tfrac12\Delta$ (key), the one-loop part of $D$ on an operator is $\gamma$, twice the map of Theorem \ref{thm:c10-brst} (block after Theorem \ref{thm:c6-dil}), and Corollary \ref{cor:c10-beisert} is $\gamma=\frac{g^2N}{8\pi^2}H$ on the scalar, $\mathfrak{su}(2)$, $\mathfrak{su}(1|1)$ and twist-two sectors. [B03c] fixes $H$ by $\mathfrak{psl}(4|4)$ invariance and a point-splitting computation in the $\mathfrak{sl}(2)$ sector; here the same numbers come from the proper-time cutoff, and their agreement is the scheme independence of the one-loop spectrum.
\end{remark}

% ANCHOR: after_label thm:c10-D1
\begin{remark}[Physics notation: $D_1$ and the anomalous dimensions]
On an operator of ghost number $0$ the physics dilatation acts through $-\Lambda^{(1)}(D)=\tfrac12\Lambda^{(1)}(\Delta)$, so its one-loop part is $\tfrac12D_1$. On an eigen-operator $\tfrac12D_1O=\gamma_OO$, with $\gamma_O$ twice the eigenvalue of the map of Theorem \ref{thm:c10-brst}, when the composite term minus the bracket of $O$ with $\hbar F_1^{(1)}(\Delta)$ beyond the field-strength renormalisation vanishes on cohomology; on the part with scalar outputs of the planar scalar sector this difference is $-(\tfrac13+\theta)$ units times $\{S^{(1)}(\Delta),O\}$, which vanishes at the value $\theta=-\tfrac13$ fixed by the arity-two equation (Theorem \ref{thm:c9-kappa}). The one-body terms of the table are twice the gauge-dependent wave-function renormalisations of the letters in Feynman gauge plus the finite counterterm of the dilatation vertex beyond the field-counting functional, and the composite term is the contribution of the line from $c$ to $A_m$ through the dilatation vertex that exists because $(Q^{GF})^2\neq0$ in Feynman gauge (block after Section 9.2). $v.X=2(x^m\partial_m+w_X)X$ is twice the classical scale transformation of the letter.
\end{remark}

% ANCHOR: after_label prop:c10-scalarsector
\begin{remark}[Physics notation: the planar scalar sector]
For $\operatorname{tr}(A_{I_1}\cdots A_{I_L})$ at large $N$, on the part with scalar outputs, $\tfrac12D_1=\frac{g^2N}{8\pi^2}\sum_l\big(1-P_{l,l+1}+\tfrac12K_{l,l+1}\big)-\tfrac12\big(\tfrac13+\theta\big)\big(-\frac{g^2N}{32\pi^2}\big)\cdot2\big(x^m\partial_m+L\big)$: the $SO(6)$ chain plus a multiple of the classical scale transformation, which vanishes at $\theta=-\tfrac13$. $P_1$ is the one-loop correction to the translation $P_m=\partial_m$ on operators; it is the multiple $-(\tfrac13+\theta_P)\big(-\frac{g^2N}{32\pi^2}\big)\partial_b$ of the classical translation, and the finite operator renormalisation $O\mapsto O+\tfrac a2\cdot2(x^m\partial_m+w_O)O$ of Proposition \ref{prop:c9-scheme}, $a$ a multiple of the unit, a position-dependent rescaling of the operator, removes it without changing the anomalous dimensions.
\end{remark}

% ANCHOR: end_of_section 10.9
\begin{remark}[Physics notation: one-loop corrections to the superconformal generators]
$Q_1$, $S_1$, $K_1$ are the one-loop corrections to the action of the supercharges $Q^i_\alpha$, $\bar Q_{\dot\alpha i}$, the special supercharges $S^\alpha{}_i$, $\bar S^{\dot\alpha i}$ and the special conformal generators on composite operators (block after Definition \ref{def:c8-components}); by Proposition \ref{prop:c10-arity1} each is a two-body plus a one-body operator on the letters, like the dilatation operator. A supercharge raises the dimension by $\tfrac12$, so its two-body kernel on a pair of total dimension $w$ and ghost number $d$ has dimension $w+\tfrac12$ and at most $w+\tfrac12+d$ output fields, and the special supercharges and special conformal generators, which carry the position $x$, lower the dimension by $\tfrac12$ and $1$ and may contain $x$ once and twice. In the canonical normalisation the terms of the supercharges of order $g$ that change the number of fields, from the commutator terms of $F_A$ in $\tfrac12(\rho(F_A)\xi,\lambda^*)$, belong to the classical $S^{(1)}(\xi)$ (block after Definition \ref{def:c2-S1}), and the one-loop components start at order $g^2N$, a term with $r$ outputs from a pair carrying $g^r$ as in the block after Proposition \ref{prop:c10-bound}. The second-order operator on pairs of letters of the holomorphic twist [BK, (3.7)] is to be compared with the two-body part of the twisted one-loop differential of Chapter 8, which contains $Q_1$ along a square-zero supercharge (block after Section 8.6).
\end{remark}
